\documentclass[a4paper,11pt]{article}

\usepackage[margin=2.5cm]{geometry}
\usepackage{listings}
\usepackage{graphicx}
\usepackage{hyperref}

\title{Practical Work 1: TCP File Transfer}
\author{Nguy\~en Tr\!ọng B\'ach \\ Student ID: 23BI14057}
\date{\today}

\begin{document}
\maketitle

\section{Goal}

The goal of this practical work is to implement a one-to-one file transfer
application over TCP/IP in the command-line interface.  The system consists of
one server and one client using the socket API in C.

\section{Protocol Design}

The protocol between client and server is very simple:

\begin{enumerate}
  \item The client establishes a TCP connection to the server.
  \item The client sends a 4--byte unsigned integer in network byte order
        representing the file size in bytes.
  \item The client sends exactly \texttt{file\_size} bytes of file data.
  \item The server receives the file size and then reads file data until all
        bytes have been received.  The server writes the data to an output file.
\end{enumerate}

This design follows the classical client/server flow shown in the lecture
slides: \texttt{socket()}, \texttt{bind()}, \texttt{listen()}, \texttt{accept()},
\texttt{connect()}, \texttt{read()}, \texttt{write()}, and \texttt{close()}.

\section{System Organisation}

The system contains two programs:

\begin{itemize}
  \item \textbf{Server} (\texttt{server.c}): listens on a TCP port, accepts one
        incoming connection and receives a file.
  \item \textbf{Client} (\texttt{client.c}): connects to the server and sends a
        file read from disk.
\end{itemize}

The compilation is managed by a simple \texttt{Makefile}.  Figure~\ref{fig:flow}
summarises the control flow.

\begin{figure}[h]
  \centering
  \begin{tabular}{p{0.45\textwidth}p{0.45\textwidth}}
    \textbf{Client} & \textbf{Server} \\
    \hline
    \texttt{socket()} & \texttt{socket()} \\
    \texttt{connect()} & \texttt{bind()}, \texttt{listen()}, \texttt{accept()} \\
    send size, send data & recv size, recv data \\
    \texttt{close()} & \texttt{close()}, \texttt{close()}
  \end{tabular}
  \caption{Client/server session for file transfer.}
  \label{fig:flow}
\end{figure}

\section{Implementation Details}

The main implementation points are:

\begin{itemize}
  \item Use of \texttt{setsockopt()} with \texttt{SO\_REUSEADDR} on the server
        to avoid ``address already in use'' errors.
  \item Use of \texttt{htonl()} and \texttt{ntohl()} to ensure the file size is
        transmitted in network byte order.
  \item A helper function \texttt{send\_all()} on the client to guarantee all
        bytes are written to the socket.
  \item A loop on the server side to receive chunks of data until the announced
        file size is fully received.
\end{itemize}

Code snippets for the critical parts of the implementation are shown in
Listings~\ref{lst:client} and \ref{lst:server}.

\subsection*{Client Snippet}
\begin{lstlisting}[language=C,caption={Sending the file size and data},label={lst:client}]
uint32_t file_size = (uint32_t)fsize;
uint32_t size_net = htonl(file_size);
send_all(sockfd, &size_net, sizeof(size_net));

while ((nread = fread(buffer, 1, BUF_SIZE, fp)) > 0) {
    send_all(sockfd, buffer, nread);
}
\end{lstlisting}

\subsection*{Server Snippet}
\begin{lstlisting}[language=C,caption={Receiving the file},label={lst:server}]
uint32_t size_net;
recv(connfd, &size_net, sizeof(size_net), MSG_WAITALL);
uint32_t file_size = ntohl(size_net);

while (received < file_size) {
    size_t to_read = (file_size - received > BUF_SIZE)
                   ? BUF_SIZE : file_size - received;
    n = recv(connfd, buffer, to_read, 0);
    fwrite(buffer, 1, n, fp);
    received += n;
}
\end{lstlisting}

\section{Work Division}

This practical was done individually.  The tasks were:

\begin{itemize}
  \item Protocol design and system organisation.
  \item Implementation of client and server in C.
  \item Testing the application with different file sizes.
  \item Writing this report in \LaTeX.
\end{itemize}

\end{document}
